\section{Representing mathematical concepts by finite observations}
\label{sec:representations}

FINER is not tied to one mathematical area.  Its generality comes from a strict rule: an ``object'' is used only through a primitive-recursive family of finite observations and explicit laws relating those observations.  This section gives reusable interfaces.

\subsection{Natural and integer algebra}

Natural-number algebra is native.  Signed integers are represented by differences \((a,b)\), denoting \(a-b\), with equality
\[
  (a,b)\sim(c,d)\quad\Longleftrightarrow\quad a+d=b+c.
\]
The existing \texttt{Core.Integers} development implements this representation and reduces polynomial identities to ring normalization in the metatheory followed by checked object derivations.  Rational numbers can be represented by signed numerators and positive denominators; equality and order are cleared by cross multiplication.  Division is never used without a positive denominator or a quantitative apartness guard.

\subsection{Finite algebraic structures}

A finite group, ring, category, graph, simplicial complex, or incidence structure is a tuple of bounds and operation/relation tables.  Its axioms are bounded formulas and therefore have reflected bits.  Theorems about one supplied finite structure reduce to finite table checks plus bounded induction.

A family with size parameter \(n\) uses primitive-recursive table accessors and a variable-length fold.  A general theorem such as associativity of a generated operation still requires an object proof; merely checking finitely many instances does not establish a universal family theorem.  Conversely, a theorem about a particular finite presentation should not be burdened with an abstract infinite structure.

\subsection{Regular rational names}

A real-like quantity is represented by a primitive-recursive rational approximation \(a(k)\) and a regularity theorem such as
\begin{equation}
  |a(k)-a(\ell)|\leq 2^{-k}
  \qquad(k\leq\ell).
  \label{eq:regular-name}
\end{equation}
All powers of two and cleared rational inequalities are primitive recursive.  The name is not a completed set-theoretic real inside \iPRA; it is a family of finite claims.  Equality at precision \(k\) means a specified rational error bound, not literal code equality.

Operations on names require precision transformers.  For addition, precision \(k+1\) for both inputs suffices for precision \(k\) of the sum.  Multiplication additionally needs magnitude bounds.  Reciprocal needs a positive apartness margin.

\begin{principle}[Guarded totalization]
Every partial operation is implemented as a total primitive-recursive program with a default output, together with a guard and a correctness theorem under that guard.  The guard must contain the numerical information needed by the error proof.
\end{principle}

For reciprocal, \(x\neq0\) is generally insufficient as finite numerical data.  A usable guard is
\begin{equation}
  |x|\geq 2^{-m}
  \label{eq:apartness-margin}
\end{equation}
for an explicit \(m\).  From approximations to \(x\) finer than \(m+k+O(1)\), rational interval arithmetic produces a \(2^{-k}\)-approximation to \(1/x\).  The demand transformer records this dependence.

\subsection{Functions and operators}

A represented function is a fixed primitive-recursive evaluator
\[
  F(k,x_{\leq d(k)})
\]
together with an error or modulus theorem.  There is no higher-type oracle in the object language; the theorem-specific evaluator reads a finite code or a finite table.  Uniform continuity on a represented compact domain is expressed by a primitive-recursive modulus \(\omega(k)\), and a finite-net theorem uses that modulus explicitly.

When an application genuinely needs an input function as a parameter, use one of three safe interfaces:
\begin{enumerate}[label=(\roman*),leftmargin=2.3em]
\item a fixed function symbol selected at the metalevel;
\item a bounded table sufficient for the current query;
\item a code for a restricted program class with a separately verified clocked evaluator.
\end{enumerate}
An unrestricted total universal primitive-recursive evaluator is unavailable by the diagonal barrier in \cref{sec:design}.

\subsection{Finite measures and probability laws}

A finite rational probability law is a list of natural weights \(w_i\) with positive total \(W\), and an observable is a rational table \(X_i\).  Expectations are represented without division by the pair
\[
  S=\sum_i w_iX_i,\qquad W=\sum_i w_i,
\]
or by cleared identities.  Linearity, monotonicity, indicators, variance identities, Markov and Chebyshev inequalities, finite conditioning, and finite independence are bounded-sum theorems.

A countable or continuous law is not imported wholesale.  It is supplied through finite quantizations \(\mu_k\), tail bounds, moment bounds, anti-concentration data, or an integration modulus.  Every limit claim is reduced to a finite query whose error budget names the required residual data.  This prevents a semantic measure-theory theorem in Lean from being mistaken for an \iPRA{} derivation.

\subsection{Compact spaces and finite nets}

A compact metric presentation consists, for each mesh parameter \(k\), of:
\begin{itemize}[leftmargin=2em]
\item a finite list of points \(N_k\);
\item a coverage certificate: every represented point lies within \(\delta_k\) of some element of \(N_k\);
\item primitive-recursive distance approximations with error bounds;
\item refinement or comparison data when cross-scale arguments need them.
\end{itemize}
Abstract compactness is replaced by the finite coverage formula actually consumed by the proof.  A minimum or finite subcover is then found by bounded search.  The production of \(N_k\) and its coverage proof may be trivial for boxes and finite complexes, or a serious residual for a general represented space.

\subsection{Topology by finite complexes and carriers}

A topological observation at precision \(k\) is a finite simplicial or cubical complex, together with carrier maps linking consecutive scales.  Chains, boundaries, cycles, homologies over a fixed finite coefficient ring, and simplicial maps are finite arrays.  A map that is not uniquely discretized may return a finite cloud; a strong carrier condition makes every selection simplicial, and contiguity yields an explicit finite chain homotopy.

This interface supports finite homology observations and quantitative separation arguments.  It does not by itself prove that a given continuous space has a suitable tower or that the tower completely captures an infinitary topological invariant.  Those are producer and comparison theorems.

\subsection{Algebraic and Diophantine geometry}

Polynomial systems are exact finite data.  Localizations, ideal-membership witnesses, Macaulay matrices, resultants, and bounded-height point searches are finite algebraic certificates.  Quotient equality should be represented by explicit witnesses, such as an ideal combination or a killing exponent, rather than by a semantic quotient type.

For rational or integral points, a verifier can check a point and can enumerate all points below a supplied height.  Completeness requires a height bound or a complete-list theorem.  Cardinality finiteness alone is not an enumeration bound; this issue is treated in \cref{sec:difficulties}.

\subsection{Observational setoids}

Code equality is usually too strict.  A represented concept should carry a reflected equivalence relation \(E\) and congruence proofs for every operation and observer used by the theorem.  This is an \emph{observational setoid}.  One works with codes throughout, transports predicates along \(E\), and uses quotient observers as in \cref{sec:observational-descent}.  There is rarely a need to construct an actual quotient sort inside the single-sorted arithmetic language.

\subsection{A uniform interface record}

The following metalevel data cover most applications:
\begin{align*}
 &\mathsf{Query},\quad \mathsf{CodeBound}(q),\quad \mathsf{Valid}(q,c),\\
 &\mathsf{Observe}(q,c),\quad \mathsf{Equiv}(q,o,o'),\\
 &\mathsf{Refine}(q,q',c,c'),\quad \mathsf{Error}(q),\\
 &\mathsf{ProducerResiduals}(q),\quad \mathsf{ConsumerGoal}(q,o).
\end{align*}
After specialization, every item is a fixed primitive-recursive symbol or formula, and every law is either an inherited/derived object sequent or an explicit residual.  This record is not meant to be one monolithic Lean type imported by every theorem; it is a checklist for designing theorem-specific interfaces.
